\section*{الفصل \ref{ch:b2:phylogenetic-trees} --- الأشجار النشوئية والساعة الجزيئية}

\begin{solution}{exo:b2:phylogenetic-trees:1}
\omterm{def:b2:phylogenetic-trees:tree}{التماثل}: صفة موروثة من سلف مشترك --- كالعضد
والكعبرة والزند في جناح خفاش وزعنفة حوت.
\omterm{def:b2:phylogenetic-trees:tree}{والتشابه العارض}: صفة متشابهة تطورت منفصلةً --- كجناحي
الخفافيش والطيور بوصفهما سطحي طيران (فالعظام متماثلة بوصفها
أطرافًا أمامية، أما الجناح بما هو جناح فلا)، أو كالعين الكاميرية عند
الأخطبوط والفقاريات. \omterm{def:b2:phylogenetic-trees:tree}{والصفة المشتقة المشتركة}: \omterm{def:b2:phylogenetic-trees:tree}{تماثل} مشتق يشترك فيه
\omterm{def:b2:phylogenetic-trees:tree}{فرع حيوي} ويحدده --- كالريش للطيور، والبيضة السلوية
للسلويات.
\end{solution}

\begin{solution}{exo:b2:phylogenetic-trees:2}
الطيور: أحادية العرق. والزواحف دون الطيور: شبه عرقية. والأسماك
دون رباعيات الأطراف: شبه عرقية (فالأسماك الرئوية أقرب إلينا منها إلى
التراوت). والفقاريات الطائرة (الخفافيش والطيور والتيروصورات): متعددة العرق.
والثدييات: أحادية العرق. وبدائيات النوى: شبه عرقية (فالبكتيريا
والعتائق، وحقيقيات النوى متروكة خارجًا --- والعتائق أقرب إلى
حقيقيات النوى منها إلى البكتيريا).
\end{solution}

\begin{solution}{exo:b2:phylogenetic-trees:3}
شقيقة الطيور: التماسيح. وشقيقة الثدييات: سائر
السلويات (السلاحف والسحالي والتماسيح والطيور معًا). وأصغر
\omterm{def:b2:phylogenetic-trees:tree}{فرع حيوي} يضم السحالي والطيور: (السحالي،(التماسيح، الطيور))،
أي مزدوجات القوس دون السلاحف على هذه الشجرة. وتبديل التماسيح
والطيور عند عقدتهما لا يغير شيئًا: فالعقدة تُدار بحرية.
\end{solution}

\begin{solution}{exo:b2:phylogenetic-trees:4}
غير المجذّرة: $3$؛ و$3\times 5\times 7 = 105$؛ و$3\times 5\times
7\times 9\times 11\times 13\times 15 = \num{2027025}$. والأشجار المجذّرة
لأربع وحدات: لكل من الأشجار الثلاث غير المجذّرة 5 فروع، فتكون 15.
\end{solution}

\begin{solution}{exo:b2:phylogenetic-trees:5}
$((A,B),(C,D))$: المواقع 1 و2 و3 خطوة واحدة لكل منها، والموقع 4 خطوتان
($\mathrm{A}$ و$\mathrm{G}$ على الجهتين)، والموقع 5 واحدة: أي 6 خطوات.
و$((A,C),(B,D))$: $2 + 2 + 2 + 1 + 1 = 8$. و$((A,D),(B,C))$: $2 + 2 + 2 +
2 + 1 = 9$. فالشجرة الأولى هي الأكثر اقتصادًا؛ وعليها يكون الموقع 4
متشابهًا عرضًا (إذ يقع التغير $\mathrm{A}\to\mathrm{G}$، أو عكسه،
مرتين).
\end{solution}

\begin{solution}{exo:b2:phylogenetic-trees:6}
صِلْ $A$ و$B$ عند الارتفاع 3. ثم $d(AB,C) = 14$ و$d(AB,D) = 16$
و$d(C,D) = 10$: فصِلْ $C$ و$D$ عند الارتفاع 5. ثم $d(AB,CD) = (14 + 16
+ 14 + 16)/4 = 15$: فالجذر عند الارتفاع $7.5$. والشجرة $((A,B),(C,D))$.
\end{solution}

\begin{solution}{exo:b2:phylogenetic-trees:7}
$d = -0.75\ln(1 - 4p/3)$: أي $0.052$ و$0.38$ و$1.21$. وعند $p = 0.6$ يكون
وسيط اللوغاريتم $0.2$ ويكون الميل $\mathrm{d}d/\mathrm{d}p
= 1/(1 - 4p/3) = 5$: فخطأ معاينة قدره $0.02$ في $p$ يصير
$0.1$ في $d$، وافتراض تساوي المعدلات في كل موقع،
وهو كاذب في كل مورّثة حقيقية، يضيف انحيازًا بالحجم نفسه. فالمورّثة
مشبعة.
\end{solution}

\begin{solution}{exo:b2:phylogenetic-trees:8}
$t = d/2k = 0.02/(2\times 10^{-9}) = 10$ ملايين سنة. وعند 500
مليون سنة يكون $d = 2\times 10^{-9}\times 5\times 10^{8} = 1.0$ ويكون
$p = 0.75(1 - e^{-4/3}) = 0.55$: أي ثلاثة أرباع الطريق إلى
الإشباع، حيث يكون التصحيح شديد الميل والتاريخ غير موثوق
--- فيلزم لمثل هذه الانفصالات مورّثة أبطأ.
\end{solution}

\begin{solution}{exo:b2:phylogenetic-trees:9}
\qty{52}{\%}: يُستعاد \omterm{def:b2:phylogenetic-trees:tree}{الفرع الحيوي} في نصف العينات بالكاد؛
فالمعطيات شبه صامتة عنه، والبدائل مجتمعةً
مرجحة مثله. و\qty{99}{\%}: الإشارة متسقة عبر
المواقع؛ فالفرع مدعوم بقوة (بمعطى النموذج
والاصطفاف --- فالأخطاء المنهجية \omterm{prop:b2:phylogenetic-trees:pitfalls}{كانجذاب الفروع الطويلة}
لا يكشفها البوتستراب). وللأول: أضف متواليات (مورّثات
أكثر)، وأضف وحدات تصنيفية تكسر الفروع، وتحقق بطريقة أخرى.
\end{solution}

\begin{solution}{exo:b2:phylogenetic-trees:10}
سلالتان تغيرت كل منهما في كسر كبير من المواقع
يكون لهما، عند أي موقع تغيرتا فيه معًا، فرصة واحدة من أربع
في الوقوع على الأساس نفسه (بأربعة أسس بمعدلات متساوية). فإن كانت كل
منهما قد تغيرت في \qty{40}{\%} من المواقع، فقد تغيرتا معًا في
\qty{16}{\%} وتتفقان بالمصادفة في \qty{4}{\%} من كل المواقع ---
وذلك قد يفوق كسر المواقع الحاملة للصفات المشتقة
المشتركة الحقيقية التي تصل كلًا منهما بأقربائها الحقيقيين
البطيئي التطور. والاقتصاد يعدّ التطابقات دون أن يسأل لماذا وقعت
فيصل الاثنتين. وإضافة وحدات تصنيفية تتفرع على امتداد كل فرع
طويل تقسمه إلى قطع أقصر، فتتوزع التطابقات
بالمصادفة على عدة فروع أقصر وتصير الصفات
المشتقة المشتركة لكل قطعة مرئية؛ والنموذج الأرجحي
الذي يتوقع تطابقات مصادفة كثيرة على الفروع الطويلة يصحح
لها مباشرةً.
\end{solution}

\begin{solution}{exo:b2:phylogenetic-trees:11}
$d = -0.75\ln(1 - 0.12) = 0.096$؛ و$t = d/2k = 0.096/(4\times 10^{-8})
= 2.4$ مليون سنة. والاستبدالات: $0.096\times \num{16000} =
1500$، والدقة النسبية $1/\sqrt{1500} = \qty{2.6}{\%}$: أي
خطأ إحصائي قدره $\pm 0.06$ مليون سنة. لكن المعدل
معاير على انفصالات أخرى وقد يكون خاطئًا بعامل اثنين
لهذه السلالة (لتباين المعدلات بين السلالات، وللفرق
بين معدل النسب المقيس على مدى أجيال
والمعدل المقيس على مدى ملايين السنين)؛ كما أن تباعد المورّثة
يسبق تباعد النوعين، إذ كانت الجمهرة السلفية
متعددة الأشكال --- فالانفصال الحقيقي أحدث
من انفصال المورّثة بزمن يتوقف على حجم الجمهرة
السلفية. (والتاريخ المقبول، من مورّثات نووية كثيرة
وأحافير، هو 6 إلى 7 ملايين سنة: فالساعة المتقدرية هنا
أسرع مما ينبغي.)
\end{solution}

\begin{solution}{exo:b2:phylogenetic-trees:12}
\omterm{prop:b2:phylogenetic-trees:pitfalls}{الفرز الناقص للسلالات}: حين يتقارب انتواعان في
الزمن، قد يُفرز تعدد شكل سلفي فتجمع شجرة المورّثة
الأنواع جمعًا مختلفًا عن شجرة النوع؛ والنقل
الأفقي: فالمورّثة المتلقاة من سلالة أخرى تحمل تاريخ تلك
السلالة؛ وتضاعف المورّثات: فمقارنة نظير في نوع
بالمثيل في آخر تعطي تاريخ التضاعف
لا تاريخ الانتواع. وتُنال شجرة النوع باستعمال
مورّثات كثيرة من أنحاء المورثوم وأخذ الشجرة التي تدعمها
الأكثرية (أو بنموذج يتوقع كسرًا معلومًا من المورّثات المتنافرة)،
وباختيار المثائل بعناية، وبتفضيل المورّثات
الريباسية والمعلوماتية في بدائيات النوى، وهي نادرًا ما تُنقل.
\end{solution}

\begin{solution}{pb:b2:phylogenetic-trees:1}
\textbf{1.} المواقع 1 و2 و4 (وهي حالات يشترك فيها $W$ و$H$ وحدهما،
مشتقة نسبةً إلى الجمل): صفات مشتقة مشتركة للمجموعة $\{W,H\}$. والموقعان 3
و6: صفتان مشتقتان مشتركتان للمجموعة $\{W,H,C\}$. والموقع 5 صفة مشتقة منفردة
للنوع $W$، أي تغير في سلالة واحدة لا يجمع شيئًا؛ ولا موقع
متشابه عرضًا على الشجرة الجزيئية.
\textbf{2.} تغير واحد لكل موقع: أي 6 خطوات.
\textbf{3.} المواقع 1 و2 و4 تحتاج الآن تغيرين لكل منها، والمواقع 3 و5
و6 تغيرًا واحدًا: أي 9 خطوات.
\textbf{4.} شجرة الحوت وفرس النهر، بثلاث خطوات. وثلاثة مواقع من
ألف هامش صغير قد تقلبه بضعة تشابهات عارضة؛
وتقويته يعني مورّثات أكثر ووحدات تصنيفية أكثر (مجترات
أخرى، والببريات، والأيائل) وبوتستراب.
\textbf{5.} يتكرر \omterm{def:b2:phylogenetic-trees:tree}{الفرع الحيوي} في خمس عينات من ست: أي دعم
متوسط، فوق المصادفة بكثير لكنه دون \qty{95}{\%}
التي تُسمى اصطلاحًا دعمًا قويًا.
\textbf{6.} على \omterm{meth:b2:phylogenetic-trees:parsimony}{المجموعة الخارجية} أن تقع خارج المجموعة لكن قريبة بما يكفي
لأن تبقى متواليتها قابلة للاصطفاف وغير مشبعة؛ والجمل
شفعي أصابع خارج مجموعة البقرة والخنزير وفرس النهر والحوت. أما
السمكة فكانت ستختلف عند مستويات شبه مشبعة، ولكانت حالاتها
غير مخبرة عن أي الحالات سلفية داخل الثدييات،
ولكان فرعها الطويل يجذب أسرع سلالات المجموعة الداخلية
تطورًا.
\textbf{7.} $d = 0.0625$ و$0.107$ و$0.131$ و$0.155$.
\textbf{8.} صِلْ $W$ و$H$ عند الارتفاع $0.031$. ثم $d(WH,C) =
0.107$ و$d(WH,P) = 0.131$ و$d(C,P) = 0.131$، وكل شيء إلى $M$
$0.155$: فصِلْ $WH$ و$C$ عند $0.054$. ثم $d(WHC,P) = 0.131$: فصِلْ
عند $0.065$. وأخيرًا صِلْ $M$ عند $0.078$. والشجرة $(M,(P,(C,(W,H))))$.
\textbf{9.} نعم: الشجرة نفسها والتعشيش نفسه.
\textbf{10.} الحوت الذي تغير أسرع مرتين كان سيصير
أبعد عن كل شيء، ومنه فرس النهر؛ وطريقة UPGMA، التي تضع كل
طرف عند الارتفاع صفر وتصل أقرب زوج، كانت ستصل فرس النهر
بالبقرة أولًا وتترك الحوت خارجًا --- أي شجرة
التشريحيين بالضبط، لسبب خاطئ.
\textbf{11.} $0.0625\times 1000 = 62.5$ استبدالًا؛ والانحراف البواسوني
المعياري $\sqrt{62.5} = 7.9$.
\textbf{12.} $7.9/62.5 = \qty{13}{\%}$.
\textbf{13.} $k = d/2t = 0.107/(1.2\times 10^{8}) = 8.9\times 10^{-10}$
في الموقع في السنة.
\textbf{14.} $t = 0.0625/(2\times 8.9\times 10^{-10}) = 35$ مليون
سنة.
\textbf{15.} الخنزير: $0.131/(1.79\times 10^{-9}) = 73$ مليون سنة؛
والجمل: $0.155/(1.79\times 10^{-9}) = 87$ مليون سنة.
\textbf{16.} $35 \pm 4.5$ مليون سنة (\qty{13}{\%}).
\textbf{17.} يتناسب المعدل مع $1/t_{\text{المعايرة}}$ ويتناسب التاريخ مع
$t_{\text{المعايرة}}$: أي $\pm 5/60 = \qty{8}{\%}$، ومن ثم $\pm 3$ ملايين سنة
--- ومجموعةً مع الخطأ البواسوني، نحو $\pm 5$ ملايين سنة.
\textbf{18.} $d = -0.75\ln(1 - 0.6) = 0.69$: فالمورّثة تغيرت في
ثلثي مواقعها فعليًا، وضرب التصحيح
$p$ في $1.5$ وصار ميله $2.5$؛ وهي قرب الإشباع لهذا
الانفصال وتاريخها قليل القيمة.
\textbf{19.} غياب صفة ليس حالة مشتقة مشتركة:
فالحيتان فقدت الطرف الخلفي وكاحله بالكلية، فلا يمكن
قول شيء عن أي كاحل كان سيكون لها. \omterm{def:b2:phylogenetic-trees:tree}{فالصفة المشتقة المشتركة}
شاهد على \omterm{def:b2:phylogenetic-trees:tree}{فرع حيوي}؛ أما فقدها في سلالة فليس شاهدًا
ضد الانتماء.
\textbf{20.} الحيتان الأحفورية لها الكاحل ذو البكرة المزدوجة: فالحيتان
تنحدر من سلف كان يملكه، وذلك يضعها داخل
شفعيات الأصابع. فالتنبؤ الجزيئي تؤكده
الصخور.
\textbf{21.} شبه عرقية: سلف مع كل خلفه
إلا خط الحوت.
\textbf{22.} فرز ناقص لتعدد شكل سلفي
عبر انتواعين سريعين (فالبقرة وفرس النهر والحوت انفصلت في غضون
بضعة ملايين من السنين)، أو مقارنة نظائر. ويُحسم الأمر
بتسلسل مورّثات كثيرة وأخذ الشجرة التي يدعمها معظمها، وتدعمها
المتسلسلة المجمّعة؛ مع التحقق من التضاعف في كل أسرة.
\textbf{23.} $d = -0.75\ln(1 - 0.333) = 0.30$؛ و$k = d/2t =
0.30/(7\times 10^{7}) = 4.3\times 10^{-9}$ في الموقع في السنة، أي خمسة
أمثال المورّثة النووية.
\textbf{24.} DNA المتقدري تنسخه بوليميراز أضعف
تدقيقًا، وهو معرض لجذور الأكسجين الحرة في سلسلة
التنفس، ولا تصلحه الآلة النووية؛ وهو
يفتقر أيضًا إلى إعادة التركيب، فيتراكم التلف.
\textbf{25.} الشجرة $(M,(P,(C,(W,H))))$: فالحيتان المجموعة الشقيقة
لفرس النهر؛ و$k = 8.9\times 10^{-10}$ في الموقع في السنة؛ وتباعد
الحوت وفرس النهر $35 \pm 5$ مليون سنة.
\end{solution}
